\documentclass[11pt]{article}
\usepackage[a4paper,margin=26mm]{geometry}
\usepackage{amsmath,amssymb,mathtools}
\usepackage{unicode-math}
\setmainfont{Latin Modern Roman}
\setmathfont{Latin Modern Math}
\usepackage{xurl,hyperref}
\hypersetup{hidelinks,pdftitle={Problem statement}}
\newcommand{\sref}[2]{\hyperlink{source:#1:#2}{[S#2]}}
\newcommand{\cataloguescope}[1]{#1}
\begin{document}
\section{Conjecture 1.19 — (The volume conjecture, [198, 296]) For any knot K, 2\texorpdfstring{$\pi$}{pi}·lim N →∞ log|JN (K)| N = v3||S3− K||, (2) where||·||denotes th…}\label{statement}

\subsection{Definitions and mathematical statement}

\paragraph{Shared notation.}$\mathbb N=\{1,2,\ldots\}$, and $\mathbb Z,\mathbb Q,\mathbb R,\mathbb C$ denote the integers, rationals, reals and complex numbers. Any other conventions are specified locally.


A knot or link is a smooth oriented embedding of one or more circles in the oriented $3$-sphere, considered up to orientation-preserving ambient isotopy. Write $O$ for the unknot, which bounds an embedded disk. Write $J_N(K)$ for the Kashaev invariant $\langle K\rangle_N$, equivalently the reduced $N$-dimensional colored Jones polynomial evaluated at $q=\exp(2\pi i/N)$, with $J_N(O)=1$. Reduction removes the unknot quantum-dimension factor before evaluation at the root of unity. For the compact knot exterior $X_K=S^3\setminus\operatorname{int}\nu K$, the simplicial volume $\|X_K,\partial X_K\|$ is the infimum of sums of absolute coefficients of real relative singular cycles representing its oriented relative fundamental class. Let $v_3$ be the volume of a regular ideal hyperbolic tetrahedron.
\paragraph{Statement recorded in the supplied source.}
For every knot $K$, does the limit exist and satisfy
\[2\pi\lim_{N\to\infty}\frac{\log|J_N(K)|}{N}
=v_3\|X_K,\partial X_K\|?\]
For a hyperbolic knot the right side is its complement's complete hyperbolic volume. The recorded conjecture applies to all knots, not only hyperbolic knots. \sref{statement}{2}
\subsection{Short English statement}
Does the exponential growth of the root-of-unity colored Jones invariant recover the knot exterior simplicial volume for every knot?
\subsection{Sources}
\begin{description}
\item[\hypertarget{source:statement:1}{S1}] ULAM. AI and the UnsolvedMath Contributors, \emph{UnsolvedMath: A Curated Collection of Open Mathematics Problems}, catalogue entry. CC BY 4.0. \url{https://huggingface.co/datasets/ulamai/UnsolvedMath}.
\item[\hypertarget{source:statement:2}{S2}] T. Ohtsuki (editor), \emph{Problems on invariants of knots and 3-manifolds}, Geometry \& Topology Monographs 4 (2002), pp. 377--572 (published 2004), item 1.19 and its surrounding definitions and remarks. \url{https://arxiv.org/abs/math/0406190v1}.
\end{description}
\label{end:statement}
\end{document}
